The kinetics simulations are important for the validation of the \gls{DNP} group parameters using the \gls{MSRE} transients, as well as investigating transients which affect the security, safety, and safeguards design decisions.
For safety design, the \gls{ULOF} and \gls{UTOP} are investigated for various reactor parameters to determine the effect of the updated \gls{DNP} group parameters.
For security design, the sabotage initiated transients are assumed to be the same as those studied for the safety design.
For safeguards design, comparison is made between various diversion scenarios with different \gls{PR} times and material quantities removed.
These diversion scenarios affect the fuel composition, which affects the effective delayed neutron fraction, which will affect the behavior of the reactor during a transient.
The updated \gls{DNP} group parameters will be used to investigate the \gls{SNM} removal which could remain undetected while using previous \gls{DNP} group paramters.

The kinetics simulations will use Moltres and the \gls{MSRE} to demonstrate the differences of the \gls{DNP} group parameters on these various scenarios and how the reactivity and temperature response can vary.
The results will provide insight into which transients are more sensitive to these updated data, as well as potential concerns.
A \gls{PRK} model can also be used for simple reactivity insertion tests with uncertainty tracking to demonstrate the differences between the \gls{DNP} group parameters in a more straightforward manner.
